% Figure: Pascal's triangle to row n = 6, with the recurrence
% C(n+1, k) = C(n, k-1) + C(n, k) drawn as a directed pair of
% in-arrows for a chosen interior cell.
\begin{figure}[htbp]
\centering
\begin{tikzpicture}[
  cell/.style={draw=engblack, circle, minimum size=0.7cm,
    inner sep=0pt, font=\small},
  highlight/.style={draw=engblack, circle, minimum size=0.7cm,
    inner sep=0pt, fill=engblue!25, font=\small\bfseries},
  arr/.style={->, thick, draw=engblack!70},
  rowlabel/.style={font=\footnotesize, anchor=east},
]

\def\dy{0.85}
\def\dx{0.95}

% Row 0
\node[cell] (n00) at (0, 0) {$1$};
\node[rowlabel] at (-4*\dx, 0) {$n = 0$};

% Row 1
\node[cell] (n10) at (-0.5*\dx, -\dy) {$1$};
\node[cell] (n11) at (0.5*\dx, -\dy) {$1$};
\node[rowlabel] at (-4*\dx, -\dy) {$n = 1$};

% Row 2
\node[cell] (n20) at (-1*\dx, -2*\dy) {$1$};
\node[cell] (n21) at (0, -2*\dy) {$2$};
\node[cell] (n22) at (1*\dx, -2*\dy) {$1$};
\node[rowlabel] at (-4*\dx, -2*\dy) {$n = 2$};

% Row 3
\node[cell] (n30) at (-1.5*\dx, -3*\dy) {$1$};
\node[cell] (n31) at (-0.5*\dx, -3*\dy) {$3$};
\node[cell] (n32) at (0.5*\dx, -3*\dy) {$3$};
\node[cell] (n33) at (1.5*\dx, -3*\dy) {$1$};
\node[rowlabel] at (-4*\dx, -3*\dy) {$n = 3$};

% Row 4
\node[cell] (n40) at (-2*\dx, -4*\dy) {$1$};
\node[cell] (n41) at (-1*\dx, -4*\dy) {$4$};
\node[cell] (n42) at (0, -4*\dy) {$6$};
\node[cell] (n43) at (1*\dx, -4*\dy) {$4$};
\node[cell] (n44) at (2*\dx, -4*\dy) {$1$};
\node[rowlabel] at (-4*\dx, -4*\dy) {$n = 4$};

% Row 5 (highlight the central entry, 10, as C(5,2))
\node[cell] (n50) at (-2.5*\dx, -5*\dy) {$1$};
\node[cell] (n51) at (-1.5*\dx, -5*\dy) {$5$};
\node[highlight] (n52) at (-0.5*\dx, -5*\dy) {$10$};
\node[cell] (n53) at (0.5*\dx, -5*\dy) {$10$};
\node[cell] (n54) at (1.5*\dx, -5*\dy) {$5$};
\node[cell] (n55) at (2.5*\dx, -5*\dy) {$1$};
\node[rowlabel] at (-4*\dx, -5*\dy) {$n = 5$};

% Row 6
\node[cell] at (-3*\dx, -6*\dy) {$1$};
\node[cell] at (-2*\dx, -6*\dy) {$6$};
\node[cell] at (-1*\dx, -6*\dy) {$15$};
\node[cell] at (0, -6*\dy) {$20$};
\node[cell] at (1*\dx, -6*\dy) {$15$};
\node[cell] at (2*\dx, -6*\dy) {$6$};
\node[cell] at (3*\dx, -6*\dy) {$1$};
\node[rowlabel] at (-4*\dx, -6*\dy) {$n = 6$};

% Recurrence arrows for the highlighted cell
\draw[arr] (n41) -- (n52);
\draw[arr] (n42) -- (n52);

% Annotation
\node[font=\small, align=left, anchor=west]
  at (2*\dx, -3.0*\dy)
  {$\binom{n+1}{k} = \binom{n}{k-1} + \binom{n}{k}$};
\node[font=\footnotesize, align=left, anchor=west]
  at (2*\dx, -3.7*\dy)
  {Pascal's recurrence:\\each entry is the sum\\of the two above.};

\end{tikzpicture}
\caption[Pascal's triangle and the binomial recurrence]{Pascal's
triangle through row $n = 6$. The arrows on the highlighted entry
$\binom{5}{2} = 10$ show the recurrence $\binom{n+1}{k} = \binom{n}{k-1}
+ \binom{n}{k}$, the identity that section 1.4 proves by induction
and that the figure makes visible. The triangle is the working table
the engineer consults when expanding $(a+b)^{n}$ for small $n$ by
the binomial theorem.}
\label{fig:vol02:ch01:pascal-triangle}
\end{figure}
